\documentclass[a4paper,12pt]{article}
\usepackage{docsTemplate}
\usepackage{eurosym}
\usepackage{array}
\newcolumntype{M}[1]{>{\centering\arraybackslash}m{#1}}

\setTitle{Piano di Qualifica}
\setAuthors{Bianca Zaghetto \\ & Mattia Oliva Medin \\ & Sara Gioia Fichera}
\setVerificators{ Mattia Oliva Medin \\ & Filippo Compagno \\ & Nenad Radulovic }
\setApprovation{ }
\setVersion{0.5}
\setType{Esterno}
\setDestination{Tullio Vardanega \\ & Riccardo Cardin \\ & Team Three Technologies}

\begin{document}
    \include{docTitlepage}

    \section*{Registro delle modifiche} {
        \begin{table}[h!]
            \rowcolors{2}{lightgray}{white}
                \begin{tabularx}{\textwidth}{|l|l|l|l|L|L|}
                \begin{tabularx}{\textwidth}{|l|l|l|l|L|L|}
                \hline
                \textbf{Vers.} & \textbf{Data} & \textbf{Autore} & \textbf{Ruolo} & \textbf{Verifica} & \textbf{Oggetto} \\
                \hline
                \textbf{0.5} & 3/01/2026 & Sara Gioia Fichera & Responsabile & Mattia Oliva Medin & Aggiunta sezione test di Sistema \\
                \hline
                \textbf{0.4} & 3/01/2026 & Bianca Zaghetto & Amministratore & Filippo Compagno & Stesura checklist \\
                \hline
                \textbf{0.3} & 2/01/2026 & Bianca Zaghetto & Amministratore & Filippo Compagno & Aggiunta sezione test di Accettazione \\
                \hline
                \textbf{0.2} & 19/12/2025 & Mattia Oliva Medin & Amministratore & Nenad Radulovic & Stesura metriche\\
                \hline
                \textbf{0.1} & 18/12/2025 & Sara Gioia Fichera & Analista & Nenad Radulovic & Struttura del documento \\
                \hline
            \end{tabularx}
        \end{table}
    }
    \newpage
    
    \tableofcontents
    \newpage

    \listoffigures
    \newpage

    \listoftables
    \newpage
    
    \section{Introduzione} {
        \subsection{Scopo del documento}
            Il presente documento ha lo scopo di definire le strategie, le metodologie e le metriche adottate per garantire la qualità del prodotto software e l'efficienza dei processi di sviluppo.

        \subsection{Glossario}
            Si rimanda al documento Glossario per chiarire i termini nel dominio del progetto che possono risultare ambigui o di difficile comprensione. I vocabili presenti nel glossario sono stati segnalati seguendo la notazione:
            \begin{center}
                parola\textsubscript{G}
            \end{center}
        \subsection{Riferimenti}
            \subsubsection{Riferimenti normativi}
                \begin{itemize}
                    \item Norme di progetto
                    \item Capitolato
                    \begin{itemize}
                        \item \url{https://www.math.unipd.it/~tullio/IS-1/2025/Progetto/C3.pdf}
                         \item[] (Ultimo accesso: 2026/01/05)
                    \end{itemize}
                    \item Regolamento del Progetto Didattico
                    \begin{itemize}
                        \item \url{https://www.math.unipd.it/~tullio/IS-1/2025/Dispense/PD1.pdf}
                         \item[] (Ultimo accesso: 2026/01/05)
                    \end{itemize}
                \end{itemize}
            \subsubsection{Riferimenti informativi}
                \begin{itemize}
                    \item Norme di Progetto
                    \item Verbali interni e esterni
                    \item Analisi dei requisiti
                    \item Slide del corso: T7 - Qualità di prodotto
                    \begin{itemize}
                        \item \url{https://www.math.unipd.it/~tullio/IS-1/2025/Dispense/T07.pdf}
                         \item[] (Ultimo accesso: 2026/01/05)
                    \end{itemize}
                    \item Slide del corso: T8 - Qualità di processo
                    \begin{itemize}
                        \item \url{https://www.math.unipd.it/~tullio/IS-1/2025/Dispense/T08.pdf}
                        \item[] (Ultimo accesso: 2026/01/05)
                    \end{itemize}
                    \item Standard ISO/IEC 12207
                        \item[] (Ultimo accesso: 2026/01/05)
                    \item Glossario
                \end{itemize}
    }

    \newpage

    \section{Obiettivi di qualità}
    La presente sezione definisce gli obiettivi di qualità del progetto, specificando per ciascuna metrica definita nel documento Norme di Progetto i range di accettazione e i valori ideali prefissati dal gruppo. La classificazione segue la distinzione tra:
    \begin{itemize}
        \item Metriche di processo
        \item Metriche di prodotto
    \end{itemize}
    
        \subsection{Qualità di processo}
            In questa sezione le metriche vengono riorganizzate seguendo questa divisione: 
            \begin{itemize}
                \item Processi Primari: Obiettivi legati alla fornitura, allo sviluppo e al controllo dei costi/tempi.
                \item Processi di Supporto: Obiettivi trasversali legati alla documentazione, alla verifica e alla gestione della configurazione.
                \item Processi Organizzativi: Obiettivi legati alla gestione delle risorse e al miglioramento continuo.
            \end{itemize}
            
            \subsubsection{Processi Primari}
                \paragraph{Fornitura} \mbox{}
                    \begin{longtable}{|M{0.15\linewidth}|M{0.35\linewidth}|M{0.2\linewidth}|M{0.2\linewidth}|}

                        \caption{Valori metriche processo di fornitura} \\
                        \hline
                        \rowcolor{gray!40}
                        \textbf{ID} & \textbf{Nome Metrica} & \textbf{Valore Accettabile} & \textbf{Valore Ideale} \\
                        \hline
                        \endhead

                        MPC01 & Metriche soddisfatte & $\ge 75\%$ & $100\%$ \\
                        \hline

                        MPC02 & Earned Value & $\ge 0$ & $\le EAC$ \\
                        \hline
                        
                        MPC03 & Planned Value & $\ge 0$ & $\le BAC$ \\
                        \hline
                        
                        MPC04 & Actual Cost & $\ge 0$ & $\le EAC$\\
                        \hline
                        
                        MPC05 & Estimated to Completion & $\ge 0$ & $\le EAC$ \\
                        \hline

                        MPC06 & Estimated at Completion & $\le BAC + 5\%$ & $\le BAC$ \\
                        \hline

                        MPC07 & Cost Variance & $\ge -5\%$ & $\ge 0\%$\\
                        \hline

                        MPC08 & Schedule Variance & &\\
                        \hline

                        MPC09 & Cost Performance Variance & & \\
                        \hline

                    \end{longtable}

                    \paragraph{Sviluppo} \mbox{}

                        \begin{longtable}{|M{0.15\linewidth}|M{0.35\linewidth}|M{0.2\linewidth}|M{0.2\linewidth}|}

                        \caption{Valori metriche processo di sviluppo} \\
                        \hline
                        \rowcolor{gray!40}
                        \textbf{ID} & \textbf{Nome Metrica} & \textbf{Valore Accettabile} & \textbf{Valore Ideale} \\
                        \hline
                        \endhead

                        MPC10 & Requirements Stability Index & &  \\
                        \hline

                    \end{longtable}


            \subsubsection{Processi di Supporto}
                \paragraph{Documentazione} \mbox{}
                    \begin{longtable}{|M{0.15\linewidth}|M{0.35\linewidth}|M{0.2\linewidth}|M{0.2\linewidth}|}

                        \caption{Valori metriche processo di documentazione} \\
                        \hline
                        \rowcolor{gray!40}
                        \textbf{ID} & \textbf{Nome Metrica} & \textbf{Valore Accettabile} & \textbf{Valore Ideale} \\
                        \hline
                        \endhead

                        MPC11 & Indice Gulpease & $\ge 60\%$ & $\ge 80\%$ \\
                        \hline

                        MPC12 & Correttezza Ortografica & $0$ & $0$ \\
                        \hline

                    \end{longtable}

                \paragraph{Gestione della qualità} \mbox{}
                    \begin{longtable}{|M{0.15\linewidth}|M{0.35\linewidth}|M{0.2\linewidth}|M{0.2\linewidth}|}

                        \caption{Valori metriche processo di gestione della qualità} \\
                        \hline
                        \rowcolor{gray!40}
                        \textbf{ID} & \textbf{Nome Metrica} & \textbf{Valore Accettabile} & \textbf{Valore Ideale} \\
                        \hline
                        \endhead

                        MPC13 & Metriche non soddisfatte & $\le 2$ & $0$ \\
                        \hline

                        MPC14 & Test Pass Rate & $-$ & $-$ \\
                        \hline

                    \end{longtable}


            \subsubsection{Processi organizzativi}
                \paragraph{Gestione} \mbox{}



        \subsection{Qualità di prodotto}
            \subsubsection{Funzionalità}
                \begin{longtable}{|M{0.15\linewidth}|M{0.35\linewidth}|M{0.2\linewidth}|M{0.2\linewidth}|}

                    \caption{Valori metriche funzionalità di prodotto} \\
                    \hline
                    \rowcolor{gray!40}
                    \textbf{ID} & \textbf{Nome Metrica} & \textbf{Valore Accettabile} & \textbf{Valore Ideale} \\
                    \hline
                    \endhead

                    MPD01 & Copertura dei requisiti obbligatori & $100\%$ & $100\%$ \\
                    \hline

                    MPD02 & Copertura dei requisiti desiderabili & $0\%$ & $100\%$ \\
                    \hline

                    MPD03 & Copertura dei requisiti opzionali & $0\%$ & $100\%$ \\
                    \hline

                \end{longtable}

            \subsubsection{Affidabilità}


            \subsubsection{Usabilità}


            \subsubsection{Efficienza}
            
\newpage
    \section{Testing}
        \subsection{Test di Unità}

        \subsection{Test di Integrazione}

        \subsection{Test di Sistema}
                   \begin{longtable}{|M{1.2cm}|M{6.3cm}|M{2.5cm}|M{1.5cm}|}
                
                \caption{Test di Sistema} \\
                \hline
                \rowcolor{gray!40}
                \textbf{ID} & \textbf{Descrizione} & \textbf{Requisito} & \textbf{Stato} \\
                \hline
                \endhead

                TS01 & Verificare che il Sistema permetta all'utente di selezionare un DIP da caricare& R1-F-Ob & NI \\
                \hline

                TS02 & Verificare che il Sistema mostri un errore se il caricamento del DIP non avviene correttamente & R2-F-Ob & NI \\
                \hline

                TS03 & Verificare che l'utente possa rimuovere un DIP dal Sistema & R3-F-Ob & NI \\
                \hline

                TS04 & Verificare che l'utente sia in grado di visualizzare il contenuto e le informazioni del DIP & R4-F-Ob & NI \\
                \hline

                TS05 & Verificare che l'utente possa visualizzare i documenti e la relativa lista degli allegati & R5-F-Ob & NI \\
                \hline

                TS06 & Verificare che l'utente possa viaulizzare ciascun allegato di ogni documento & R6-F-Ob & NI \\
                \hline

                TS07 & Verificare che il sistema possa mostrare le informazioni di conservazione di un DIP & R7-F-Ob & NI \\
                \hline

                TS08 & Verificare che l'utente sia in grado di selezionare uno o più documenti contenuti nel DIP & R8-F-Ob & NI \\
                \hline

                TS09 & Verificare che l'utente sia in grado di deselezionare uno o più documenti precedentemente selezionati & R9-F-Ob & NI \\
                \hline

                TS10 & Verificare che l'utente possa selezionare "Impronta crittografica" come criterio di ricerca & R10-F-De & NI \\
                \hline

                TS11 & Verificare che l'utente possa selezionare "Identificativo documento" come criterio di ricerca & R11-F-De & NI \\
                \hline

                TS12 & Verificare che l'utente possa selezionare "Modalità di formazione" come criterio di ricerca & R12-F-De & NI \\
                \hline

                TS13 & Verificare che l'utente possa selezionare "Tipologia documentale" come criterio di ricerca & R13-F-Ob & NI \\
                \hline

                TS14 & Verificare che l'utente possa selezionare "Tipologia di flusso" come criterio di ricerca & R14-F-De & NI \\
                \hline

                TS15 & Verificare che l'utente possa selezionare "Tipo registro" come criterio di ricerca & R15-F-De & NI \\
                \hline
               
                TS16 & Verificare che l'utente possa selezionare "Data registrazione" come criterio di ricerca & R16-F-Ob & NI \\
                \hline
                
                TS17 & Verificare che l'utente possa selezionare "Ora registrazione" come criterio di ricerca & R17-F-De & NI \\
                \hline
                
                TS18 & Verificare che l'utente possa selezionare "Numero documento" come criterio di ricerca & R18-F-Ob & NI \\
                \hline

                TS19 & Verificare che l'utente possa selezionare "Codice registro" come criterio di ricerca & R19-F-Ob & NI \\
                \hline
                
                TS20 & Verificare che l'utente possa selezionare "Ruolo" come criterio di ricerca & R20-F-De & NI \\
                \hline
                
                TS21 & Verificare che l'utente possa selezionare "Tipo soggetto" come criterio di ricerca & R21-F-De & NI \\
                \hline
                
                TS22 & Verificare che l'utente possa selezionare "Cognome" come criterio di ricerca & R22-F-De & NI \\
                \hline
                
                TS23 & Verificare che l'utente possa selezionare "Nome" come criterio di ricerca & R23-F-De & NI \\
                \hline
                
                TS24 & Verificare che l'utente possa selezionare "Codice Fiscale" come criterio di ricerca & R24-F-De & NI \\
                \hline
                
                TS25 & Verificare che l'utente possa selezionare "Indirizzo digitale" come criterio di ricerca & R25-F-De & NI \\
                \hline
                
                TS26 & Verificare che l'utente possa selezionare "Denominazione Organizzazione" come criterio di ricerca & R26-F-De & NI \\
                \hline
                
                TS27 & Verificare che l'utente possa selezionare "Partita IVA" come criterio di ricerca & R27-F-De & NI \\
                \hline

                TS28 & Verificare che l'utente possa selezionare "Denominazione ufficio" come criterio di ricerca & R28-F-De & NI \\
                \hline
                
                TS29 & Verificare che l'utente possa selezionare "Denominazione amministrazione" come criterio di ricerca & R29-F-De & NI \\
                \hline

                TS30 & Verificare che l'utente possa selezionare "Codice IPA" come criterio di ricerca & R30-F-De & NI \\
                \hline

                TS31 & Verificare che l'utente possa selezionare "Denominazione amministrazione AOO" come criterio di ricerca & R31-F-De & NI \\
                \hline

                TS32 & Verificare che l'utente possa selezionare "Codice IPA AOO" come criterio di ricerca & R32-F-De & NI \\
                \hline

                TS33 & Verificare che l'utente possa selezionare "Denominazione amministrazione UOR" come criterio di ricerca & R33-F-De & NI \\
                \hline

                TS34 & Verificare che l'utente possa selezionare "Codice IPA UOR" come criterio di ricerca & R34-F-De & NI \\
                \hline
                
                TS35 & Verificare che l'utente possa selezionare "Denominazione sistema" come criterio di ricerca & R35-F-De & NI \\
                \hline
                
                TS36 & Verificare che l'utente possa selezionare "Oggetto" come criterio di ricerca & R36-F-Ob & NI \\
                \hline
                
                TS37 & Verificare che l'utente possa selezionare "Parola chiave" come criterio di ricerca & R37-F-De & NI \\
                \hline
                
                TS38 & Verificare che l'utente possa selezionare "Numero allegati" come criterio di ricerca & R38-F-De & NI \\
                \hline
                
                TS39 & Verificare che l'utente possa selezionare "Impronta crittografica allegato" come criterio di ricerca & R39-F-De & NI \\
                \hline
                
                TS40 & Verificare che l'utente possa selezionare "Identificativo allegato" come criterio di ricerca & R40-F-De & NI \\
                \hline
                
                TS41 & Verificare che l'utente possa selezionare "Descrizione allegato" come criterio di ricerca & R41-F-De & NI \\
                \hline
                
                TS42 & Verificare che l'utente possa selezionare "Indice di classificazione" come criterio di ricerca & R42-F-De & NI \\
                \hline
                
                TS43 & Verificare che l'utente possa selezionare "Descrizione dell'indice di classificazione" come criterio di ricerca & R43-F-De & NI \\
                \hline

                TS44 & Verificare che l'utente possa selezionare "Piano di classificazione" come criterio di ricerca & R44-F-De & NI \\
                \hline
            
                TS45 & Verificare che l'utente possa selezionare "Riservato" come criterio di ricerca & R45-F-De & NI \\
                \hline
                
                TS46 & Verificare che l'utente possa selezionare "Formato" come criterio di ricerca & R46-F-Ob & NI \\
                \hline
                
                TS47 & Verificare che l'utente possa selezionare "Nome prodotto software" come criterio di ricerca & R47-F-De & NI \\
                \hline
                                
                TS48 & Verificare che l'utente possa selezionare "Versione prodotto software" come criterio di ricerca & R48-F-De & NI \\
                \hline
                                
                TS49 & Verificare che l'utente possa selezionare "Produttore software" come criterio di ricerca & R49-F-De & NI \\
                \hline
                                
                TS50 & Verificare che l'utente possa selezionare "Firmato digitalmente" come criterio di ricerca & R50-F-De & NI \\
                \hline
                                
                TS51 & Verificare che l'utente possa selezionare "Sigillato elettronicamente" come criterio di ricerca & R51-F-De & NI \\
                \hline
                                
                TS52 & Verificare che l'utente possa selezionare "Marcatura temporale" come criterio di ricerca & R52-F-De & NI \\
                \hline
                                
                TS53 & Verificare che l'utente possa selezionare "Conformità coppie immagine su supporto informatico" come criterio di ricerca & R53-F-De & NI \\
                \hline
                                
                TS54 & Verificare che l'utente possa selezionare "Tipo aggregazione" come criterio di ricerca & R54-F-Ob & NI \\
                \hline
                                
                TS55 & Verificare che l'utente possa selezionare "Identificativo aggregazione" come criterio di ricerca & R55-F-De & NI \\
                \hline
                                
                TS56 & Verificare che l'utente possa selezionare "Tipologia fascicolo" come criterio di ricerca & R56-F-De & NI \\
                \hline
                                
                TS57 & Verificare che l'utente possa selezionare "Nome del documento" come criterio di ricerca & R57-F-Ob & NI \\
                \hline
                                
                TS58 & Verificare che l'utente possa selezionare "Versione del documento" come criterio di ricerca & R58-F-De & NI \\
                \hline
                                
                TS59 & Verificare che l'utente possa selezionare "Tipo modifica" come criterio di ricerca & R59-F-De & NI \\
                \hline
                                
                TS60 & Verificare che l'utente possa selezionare "Data modifica" come criterio di ricerca & R60-F-De & NI \\
                \hline
                                
                TS61 & Verificare che l'utente possa selezionare "Ora modifica" come criterio di ricerca & R61-F-De & NI \\
                \hline
                                
                TS62 & Verificare che l'utente possa selezionare "Tempo di conservazione" come criterio di ricerca & R62-F-De & NI \\
                \hline
                                
                TS63 & Verificare che l'utente possa selezionare "Note" come criterio di ricerca & R63-F-De & NI \\
                \hline
                                
                TS64 & Verificare che il Sistema permetta di effettuare una ricerca impostando uno o più criteri di ricerca & R64-F-Ob & NI \\
                \hline
                                
                TS65 & Verificare che il Sistema mostri come risultato della ricerca i documenti che soddisfano i criteri di ricerca impostati& R65-F-Ob & NI \\
                \hline
                                
                TS66 & Verificare che l'utente possa visualizzare le informazioni relative ai documenti risultanti dalla ricerca & R66-F-Ob & NI \\
                \hline
                                
                TS67 & Verificare che il sistema mostri un messaggio informativo se il risultato della ricerca è vuoto & R67-F-Ob & NI \\
                \hline
                                
                TS68 & Verificare che l'utente possa visualizzare l'anteprima di un documento selezionato & R68-F-Ob & NI \\
                \hline
                                
                TS69 & Verificare che il sistema mostri un messaggio di avviso se non è possibile visualizzare l'anteprima di un documento selezionato & R69-F-Ob & NI \\
                \hline
                                
                TS70 & Verificare che l'utente possa esportare un documento selezionato, specificando il percorso dove salvarlo & R70-F-Ob & NI \\
                \hline
                                
                TS71 & Verificare che il sistema mostri un messaggio di errore se l'esportazione fallisce & R7c-F-Ob & NI \\
                \hline
                                
                TS72 & Verificare che l'utente possa configurare la ricerca semantica inserendo una query in linguaggio naturale & R72-F-Op & NI \\
                \hline
                                
                TS73 & Verificare che il sistema mostri i risultati della ricerca semantica configurata & R73-F-Op & NI \\
                \hline
                                
                TS74 & Verificare che il sistema permetta la stampa di un documento selzionato & R74-F-Op & NI \\
                \hline
            \end{longtable}


        \subsection{Test di Accettazione}
        \begin{longtable}{|M{1.5cm}|M{8.5cm}|M{2cm}|}

                \caption{Test di Accettazione} \\
                \hline
                \rowcolor{gray!40}
                \textbf{ID} & \textbf{Descrizione} & \textbf{Stato} \\
                \hline
                \endhead

                TA01 & Verificare che l'utente possa caricare un pacchetto e che in seguito possa visionare la dashboard principale & NI \\
                \hline
                
                TA02 & Verificare che l'utente possa chiudere un pacchetto e che il sistema torni alla schermata di caricamento & NI \\
                \hline

                TA03 & Verificare che l’utente possa scorrere la lista dei documenti contenuti nel pacchetto e che possa accedere ai dettagli di un singolo documento & NI \\
                \hline
                
                TA04 & Verificare che sia possibile selezionare il criterio di ricerca “Impronta crittografica” e che venga effettuata la ricerca correttamente & NI \\
                \hline

                TA05 & Verificare che, dopo la configurazione di uno o più criteri, venga mostrata all'utente la lista di documenti del pacchetto filtrata, e nel caso di ricerca senza corrispondenze venga mostrato un messaggio di avviso & NI \\
                \hline
                
                TA06 & Verificare che il prodotto dia la possibilità di effettuare una ricerca semantica tramite una query in linguaggio naturale & NI \\
                \hline

                TA07 & Verificare che venga mostrata all’utente l’anteprima di un documento dopo aver selezionato un documento e la funzionalità di anteprima & NI \\
                \hline
                
                TA08 & Verificare che l’utente possa esportare con successo un documento a seguito della selezione e dell’inserimento del percorso di esportazione & NI \\
                \hline

                TA09 & Verificare che il prodotto dia la possibilità all’utente di stampare un documento & NI \\
                \hline

                TA10 & Verificare che l’utente possa selezionare un documento firmato digitalmente e avviare la procedura di verifica, la quale, al termine dell’elaborazione, restituisce l’esito della validazione & NI \\
                \hline

        \end{longtable}

        


        \subsection{Checklist}
            \subsubsection{Documentazione}
                \begin{longtable}{|M{3.5cm}|M{8.5cm}|}

                    \caption{Checklist documentazione} \\
                    \hline
                    \rowcolor{gray!40}
                    \textbf{Nome} & \textbf{Descrizione} \\
                    \hline
                    \endhead

                    Immagini e tabelle & Tutte le immagini e le tabelle devono essere numerate e corredate da un titolo \\
                    \hline

                    Riferimento a risorse online & Dato che le risorse online mutano continuamente è opportuno indicare accanto al link della risorsa la data di ultima consultazione \\
                    \hline

                    Termine del glossario & Ogni volta che in un documento prodotto dal gruppo compare un termine descritto nel glossario deve essere segnalato con l’apposito stile deciso, parola\textsubscript{G}.\\
                    \hline

                    Formato delle date & Tutte le date che compaiono nel nome di un documento o nella tabella di versionamento devono avere il formato AAAA-MM-GG \\
                    \hline

                    Sezioni vuote & Tutte le sezioni dei documenti prodotti devono essere compilate \\
                    \hline

                    Cambiamento di versione & Ad ogni modifica al documento prodotta deve corrispondere un aggiornamento di versione e una riga nella tabella di versionamento \\
                    \hline

                    Errori di sintassi & Non devono essere presenti errori di sintassi \\
                    \hline

                    Forma impersonale & La forma impersonale dei verbi nei verbali è da preferire \\
                    \hline

                    Acronimi in maiuscolo & Gli acronimi devono essere in maiuscolo (es. DIP) \\
                    \hline

                    Formattazione titoli & I titoli dei documenti devono avere la prima lettera maiuscola \\
                    \hline

                    Denominazione titoli dei verbali & Il titolo dei Verbali deve essere così composto: Verbale Tipologia (Interno/Esterno) Data (nel formato DD/MM/AAAA) \\
                    \hline

                    Ordine dei nomi dei componenti & I nomi dei componenti del gruppo, all’interno dei documenti prodotti, devono essere riportati in ordine alfabetico \\
                    \hline

                    Ruoli in maiuscolo & Tutti i nomi dei ruoli devono avere la prima lettera maiuscola \\
                    \hline

                    Decisioni dei verbali & Tutte le decisioni riportate nell’apposita sezione dei verbali devono essere corredate da un identificativo univoco \\
                    \hline

                \end{longtable}

            \subsubsection{Analisi dei Requisiti}

                \begin{longtable}{|M{3.5cm}|M{8.5cm}|}
                    
                    \caption{Checklist Analisi dei requisiti} \\
                    \hline
                    \rowcolor{gray!40}
                    \textbf{Nome} & \textbf{Descrizione} \\
                    \hline
                    \endhead

                    Relazione tra casi d’uso e requisiti & Per ogni requisito deve corrispondere almeno un caso d’uso \\
                    \hline

                    Numerazione dei casi d’uso & I casi d’uso devono essere numerati in ordine crescente e nel caso di sotto-casi viene preso il numero del caso d’uso principale e viene aggiunta accanto un’altra cifra partendo da 1 \\
                    \hline

                    Diagrammi UML & Estensioni, inclusioni e generalizzazioni di un caso d’uso devono essere rappresentati in unico diagramma \\
                    \hline

                \end{longtable}

        


    \section{Cruscotto}
        \subsection{M1}

        \subsection{M2}

\end{document}